\chapter{Interactions, the Heisenberg picture and Lieb--Robinson bounds}
\label{ch:dynamics}

This chapter treats the evolution generated by a family of local terms on a
system of sites, with the notation of Chapter~\ref{ch:sites}. It proves the
local commutator recursion of Hastings and Koma~\cite{HastingsKoma2006SpectralGap}
(see also~\cite{NachtergaeleOgataSims2006Propagation}), exact locality on regions
closed under the interaction, the Lieb--Robinson bound in the distance of a
graph on the sites, and the approximation of the evolved operator by its
expectation onto a graph ball. The constants depend only on the range, the
size of the supports and the norms of the terms, never on the number of sites
or the local dimensions.

\section{Interactions and the Heisenberg picture}
\label{sec:dynamics_interaction}

\begin{definition}[Interaction]
    \label{def:dynamics_interaction}
    \lean{ManyBody.Interaction, ManyBody.Interaction.hamiltonian,
        ManyBody.Interaction.neg_term, ManyBody.Interaction.hamiltonian_neg,
        ManyBody.Interaction.isHermitian_hamiltonian}
    \leanok
    \uses{thm:sites_supported_subalgebra}
    Let $I$ be a set and let $T=(T_i)_{i\in I}$ be regions. An interaction with
    supports $T$ is a family $\Phi=(h_i)_{i\in I}$ of operators with
    $h_i\in\mathcal A_{T_i}$. Several terms may share a support, and no term is
    required to be Hermitian or bounded by a given constant. When $I$ is finite,
    the Hamiltonian of $\Phi$ is $H_\Phi=\sum_{i\in I}h_i$; it is Hermitian when
    every term is. The negated interaction $-\Phi=(-h_i)_i$ has the same
    supports and the Hamiltonian $-H_\Phi$.
\end{definition}

\begin{definition}[Heisenberg picture]
    \label{def:dynamics_heisenberg}
    \lean{ManyBody.heisenbergEvolution, ManyBody.heisenbergEvolution_def}
    \leanok
    For a square matrix $H$ and $t\in\R$, the Heisenberg evolution of an
    operator $A$ is $\tau_t(A)=e^{itH}Ae^{-itH}$.
\end{definition}

\begin{lemma}[Rescaling energy and time]
    \label{lem:dynamics_time_scaling}
    \lean{ManyBody.heisenbergEvolution_smul,
        ManyBody.heisenbergEvolution_inv_smul_mul}
    \leanok
    \uses{def:dynamics_heisenberg}
    For real $c$, the evolution under $cH$ at time $t$ is the evolution under
    $H$ at time $ct$. In particular, for $\Delta\neq0$ the evolution under
    $\Delta^{-1}H$ at time $\Delta t$ is the evolution under $H$ at time $t$.
\end{lemma}
\begin{proof}\leanok
    Both sides are $e^{ictH}Ae^{-ictH}$.
\end{proof}

\section{The local commutator recursion}
\label{sec:dynamics_recursion}

\begin{lemma}[Commutator estimates in a normed ring]
    \label{lem:dynamics_commutator_ring}
    \lean{norm_commutator_le, norm_exp_commutator_eq}
    \leanok
    In a normed ring, $\lVert XA-AX\rVert\leq2\lVert A\rVert\lVert X\rVert$. In a
    $C^*$-algebra, for skew-adjoint $K$,
    $\lVert[e^{tK}Ae^{-tK},B]\rVert=\lVert[A,e^{-tK}Be^{tK}]\rVert$.
\end{lemma}
\begin{proof}\leanok
    The first is the triangle inequality. For the second, conjugation by the
    unitary $e^{tK}$ is an isometric automorphism.
\end{proof}

\begin{definition}[Restricted commutator norm]
    \label{def:dynamics_commutator_norm}
    \lean{ManyBody.heisenbergCommutatorNorm,
        ManyBody.norm_heisenberg_commutator_le_heisenbergCommutatorNorm}
    \leanok
    \uses{def:dynamics_heisenberg, thm:sites_supported_subalgebra}
    For operators $H$ and $B$, a region $S$ and $t\in\R$, let $F_S(t)$ be the
    operator norm of the linear map $A\mapsto[\tau_t(A),B]$ on $\mathcal A_S$.
    Thus $\lVert[\tau_t(A),B]\rVert\leq F_S(t)\lVert A\rVert$ for
    $A\in\mathcal A_S$.
\end{definition}

\begin{lemma}[Elementary properties of the restricted commutator norm]
    \label{lem:dynamics_commutator_norm}
    \lean{ManyBody.heisenbergCommutatorNorm_nonneg,
        ManyBody.heisenbergCommutatorNorm_neg_time,
        ManyBody.continuous_heisenbergCommutatorNorm,
        ManyBody.heisenbergCommutatorNorm_zero_le,
        ManyBody.heisenbergCommutatorNorm_zero_of_disjoint}
    \leanok
    \uses{def:dynamics_commutator_norm}
    The function $F_S$ is nonnegative and continuous. Reversing time amounts to
    replacing $H$ by $-H$. At time zero, $F_S(0)\leq2\lVert B\rVert$, and
    $F_S(0)=0$ when $B$ is supported on a region disjoint from $S$.
\end{lemma}
\begin{proof}\leanok
    \uses{lem:dynamics_commutator_ring, lem:sites_supported_commute}
    $F_S(t)$ is the norm of a linear map that depends continuously on $t$,
    through $e^{\pm itH}$. Replacing $t$ by $-t$ in $e^{itH}$ is replacing $H$
    by $-H$. At $t=0$ the map is $A\mapsto[A,B]$, whose norm is at most
    $2\lVert B\rVert$ by Lemma~\ref{lem:dynamics_commutator_ring}, and which
    vanishes on $\mathcal A_S$ when $B$ is supported on a region disjoint from
    $S$, by Lemma~\ref{lem:sites_supported_commute}.
\end{proof}

\begin{theorem}[Local commutator recursion]
    \label{thm:dynamics_recursion}
    \lean{ManyBody.heisenbergCommutatorNorm_le_integral}
    \leanok
    \uses{def:dynamics_interaction, def:dynamics_commutator_norm}
    Let $\Phi=(h_j)_{j\in J}$ be an interaction with finitely many Hermitian
    terms, supports $T_j$ and Hamiltonian $H$. For every operator $B$, region
    $S$ and $t\geq0$,
    \begin{align*}
        F_S(t)\leq F_S(0)+2\sum_{j:\,T_j\cap S\neq\emptyset}
            \lVert h_j\rVert\int_0^tF_{T_j}(s)\,ds.
    \end{align*}
\end{theorem}
\begin{proof}\leanok
    \uses{lem:dynamics_commutator_ring, lem:sites_supported_commute}
    Write $H=H_S+H'$, where $H'$ collects the terms with supports disjoint from
    $S$; by Lemma~\ref{lem:sites_supported_commute}, $H'$ commutes with every
    $A\in\mathcal A_S$. In the interaction picture with respect to $H'$,
    differentiating
    $s\mapsto e^{isH'}[A,e^{-isH}Be^{isH}]e^{-isH'}$ leaves only the
    commutator with $H_S$, and integrating from $0$ to $t$ gives
    $\lVert[A,\tau_{-t}(B)]\rVert\leq\lVert[A,B]\rVert+2\lVert A\rVert
    \int_0^t\lVert[H_S,\tau_{-s}(B)]\rVert\,ds$. Expanding $H_S$ into its terms
    and bounding each term by $\lVert h_j\rVert F_{T_j}(s)$ proves the claim;
    this is the argument of~\cite[Appendix~A]{HastingsKoma2006SpectralGap}.
    % Hastings--Koma, arXiv:math-ph/0507008, Appendix A, (A.12)--(A.14).
\end{proof}

\begin{theorem}[Exact locality on interaction-closed regions]
    \label{thm:dynamics_closed_region}
    \lean{ManyBody.commute_heisenbergEvolution_of_closed}
    \leanok
    \uses{def:dynamics_interaction, def:dynamics_heisenberg}
    Let $\Phi$ be an interaction with finitely many terms, and let $C$ be a
    region such that every support lies in $C$ or is disjoint from $C$. If
    $A\in\mathcal A_C$ and $B\in\mathcal A_D$ with $D\cap C=\emptyset$, then
    $\tau_t(A)$ commutes with $B$ for every $t$.
\end{theorem}
\begin{proof}\leanok
    \uses{lem:sites_supported_commute, lem:sites_supported_mono}
    The Hamiltonian splits as $H_C+H_O$, where $H_C\in\mathcal A_C$ collects the
    terms supported in $C$ and $H_O$ commutes with $\mathcal A_C$. Hence
    $e^{itH}=e^{itH_C}e^{itH_O}$, so $\tau_t(A)=e^{itH_C}Ae^{-itH_C}$, which lies
    in $\mathcal A_C$ and commutes with $B$.
\end{proof}

\section{Propagation in graph distance}
\label{sec:dynamics_graph}

Throughout this section $G$ is a graph on the sites with extended path
distance $d_G$, and $\Phi=(h_k)_{k\in J}$ is an interaction with finitely many
Hermitian terms, supports $T_k$ and anchors $a_k\in T_k$. We assume that every
support has $d_G$-diameter at most $R$ and at most $v_R$ sites, and that
$\sum_{k:\,x\in T_k}\lVert h_k\rVert\leq b_0$ for every site $x$.

\begin{lemma}[Weighted Volterra comparison]
    \label{lem:dynamics_volterra}
    \lean{le_mul_exp_of_le_const_add_integral, finite_volterra_le_weight_mul_exp,
        finite_volterra_zero_le_weight_mul_exp_sub_one,
        finite_rectangular_volterra_zero_le_weight_mul_exp_sub_one}
    \leanok
    Let $f_i$, $i$ in a finite set, be continuous nonnegative functions on
    $[0,T]$ with $f_i(t)\leq b_i+\sum_ja_{ij}\int_0^tf_j(s)\,ds$, where
    $a_{ij}\geq0$. If $w_i>0$, $\sum_ja_{ij}w_j\leq v\,w_i$ and $b_i\leq c\,w_i$,
    then $f_i(t)\leq c\,w_ie^{vt}$ on $[0,T]$; if moreover $b_i=0$ and $v>0$,
    then $f_i(t)\leq c\,w_i(e^{vt}-1)$. The last bound holds also when the
    generators carry their own index set: if
    $f_i(t)\leq b_i+\sum_{k}a_{ik}\int_0^tf_{e(k)}(s)\,ds$ for a map $e$ from
    the generators to the indices, with $\sum_ka_{ik}w_{e(k)}\leq v\,w_i$, so
    that generators with the same index are counted separately.
\end{lemma}
\begin{proof}\leanok
    The maximum of $f_i/w_i$ satisfies the scalar inequality
    $M(t)\leq c+v\int_0^tM$, to which Gr\"onwall's inequality applies. For
    $b_i=0$, integrate the first bound once more. With a separate index set for
    the generators, summing $a_{ik}$ over the generators $k$ with $e(k)=j$
    reduces the system to the first form.
\end{proof}

\begin{theorem}[Lieb--Robinson bound in graph distance]
    \label{thm:dynamics_lieb_robinson}
    \lean{ManyBody.heisenbergCommutatorNorm_le_graphDistance,
        ManyBody.norm_heisenberg_commutator_le_graphDistance,
        ManyBody.heisenbergCommutatorNorm_le_graphDistance_finset,
        ManyBody.norm_heisenberg_commutator_le_graphDistance_finset,
        ManyBody.heisenbergCommutatorNorm_le_graphDistance_finsets_numeric,
        ManyBody.norm_heisenberg_commutator_le_graphDistance_finsets_numeric}
    \leanok
    \uses{def:dynamics_interaction, def:dynamics_heisenberg}
    Let $\mu\geq0$, $A\in\mathcal A_{T_i}$, and let $B$ be supported on a site $y$
    with $n\leq d_G(a_i,y)$. Then for every $t\in\R$
    \begin{align*}
        \lVert[\tau_t(A),B]\rVert\leq2e^{\mu R}\lVert A\rVert\lVert B\rVert
            \exp\bigl(2v_Rb_0e^{2\mu R}|t|-\mu n\bigr).
    \end{align*}
    For an arbitrary region $S$ with $A\in\mathcal A_S$ and $d_G(x,y)\geq n$ for
    all $x\in S$, the same holds with the prefactor
    $(2+4|b_0|e^{2\mu R})|S|$ and the velocity $2v_R|b_0|e^{2\mu R}+1$. If
    $v_R\geq1$ and $B$ is supported on a region $Y$ with $d_G(x,y)\geq n$ for
    all $x\in S$ and $y\in Y$, then
    $\lVert[\tau_t(A),B]\rVert\leq4|S|\lVert A\rVert\lVert B\rVert
    \exp(4v_R|b_0|e^{2\mu R}|t|-\mu n)$, with no factor depending on $Y$.
    The first bound is~\cite[Lemma~4.1]{OpenAI2026AreaLaw} on a general finite
    graph, and the bound between two regions is the estimate used
    in~\cite{Hastings2007AreaLaw}.
\end{theorem}
\begin{proof}\leanok
    \uses{thm:dynamics_recursion, lem:dynamics_volterra, lem:dynamics_commutator_norm}
    For $t\geq0$ apply Theorem~\ref{thm:dynamics_recursion} to every support
    $T_k$. Supports that meet have anchors at distance at most $2R$, so the
    weights $w_k=e^{-\mu\min(n,d_G(a_k,Y))}$ satisfy the row bound of
    Lemma~\ref{lem:dynamics_volterra} with $v=2v_Rb_0e^{2\mu R}$, and the
    initial values are at most $2\lVert B\rVert e^{\mu R}w_k$. For $t<0$ apply
    the result to $-\Phi$. For a general region $S$, apply the bound to the
    supports meeting $S$ in the recursion for $F_S$ and absorb the factor of
    time into the velocity.
\end{proof}

\begin{theorem}[Exact vanishing at infinite distance]
    \label{thm:dynamics_infinite_distance}
    \lean{ManyBody.heisenberg_commutator_eq_zero_of_forall_edist_eq_top,
        ManyBody.heisenberg_commutator_eq_zero_of_edist_eq_top}
    \leanok
    \uses{def:dynamics_interaction, def:dynamics_heisenberg}
    If every support has finite diameter, $A\in\mathcal A_S$, $B$ is supported
    on a site $y$, and no site of $S$ is joined to $y$ by a path in $G$, then
    $[\tau_t(A),B]=0$ for all $t$.
\end{theorem}
\begin{proof}\leanok
    \uses{thm:dynamics_closed_region}
    The sites not joined to $y$ form a region closed under the interaction that
    contains $S$ and not $y$; apply Theorem~\ref{thm:dynamics_closed_region}.
\end{proof}

\begin{corollary}[Constants fixed in advance]
    \label{cor:dynamics_lieb_robinson_constants}
    \lean{ManyBody.exists_graph_lieb_robinson}
    \leanok
    \uses{def:dynamics_interaction, def:dynamics_heisenberg}
    For given $R$, $v_R$ and $b_0$ there are $C,v,c>0$ such that for every
    finite graph, every system of sites and every interaction as above,
    $\lVert[\tau_t(A),B]\rVert\leq C\lVert A\rVert\lVert B\rVert e^{v|t|-cn}$
    whenever $A\in\mathcal A_{T_i}$, $B$ is supported on $y$ and
    $n\leq d_G(a_i,y)$, and the commutator vanishes when $d_G(a_i,y)=\infty$.
    This is the form of~\cite[Lemma~4.1]{OpenAI2026AreaLaw}.
\end{corollary}
\begin{proof}\leanok
    \uses{thm:dynamics_lieb_robinson, thm:dynamics_infinite_distance}
    Take $\mu=1$ in Theorem~\ref{thm:dynamics_lieb_robinson}.
\end{proof}

\section{Localization onto graph balls}
\label{sec:dynamics_localization}

\begin{definition}[Graph ball]
    \label{def:dynamics_graph_ball}
    \lean{ManyBody.graphBall, ManyBody.coe_graphBall, ManyBody.mem_graphBall}
    \leanok
    The graph ball $N_l(a)=\{x:d_G(a,x)\leq l\}$ of natural radius $l$ is the
    open ball of radius $l+1$ in the extended distance, as a region.
\end{definition}

\begin{lemma}[Monotonicity and stabilization of graph balls]
    \label{lem:dynamics_graph_ball}
    \lean{ManyBody.graphBall_mono, ManyBody.notMem_graphBall_of_edist_eq_top,
        ManyBody.coe_graphBall_eq_reachable_of_sup_dist_le,
        ManyBody.exists_graphBall_eq_reachable}
    \leanok
    \uses{def:dynamics_graph_ball}
    The balls $N_l(a)$ increase with $l$ and contain no site at infinite
    distance from $a$. For every $l$ at least the largest finite distance from
    $a$, the ball $N_l(a)$ is the set of sites joined to $a$ by a path.
\end{lemma}
\begin{proof}\leanok
    Monotonicity and the exclusion of infinite distances are immediate from
    $d_G(a,x)\leq l$. A site joined to $a$ has finite distance, at most the
    largest finite distance from $a$, hence at most $l$; conversely a site at
    finite distance is joined to $a$.
\end{proof}

\begin{theorem}[Localization onto graph balls]
    \label{thm:dynamics_localization}
    \lean{ManyBody.sum_pow_mul_exp_neg_le,
        ManyBody.sum_exp_neg_edist_compl_graphBall_le,
        ManyBody.norm_heisenberg_sub_regionExpectation_graphBall_le,
        ManyBody.norm_heisenberg_sub_regionExpectation_graphNeighborhood_le}
    \leanok
    \uses{def:dynamics_graph_ball, def:sites_region_expectation,
        def:dynamics_interaction, def:dynamics_heisenberg}
    Assume, in addition to the hypotheses of
    Section~\ref{sec:dynamics_graph}, that at most $K_g(d+1)^k$ sites lie at
    distance exactly $d$ from any site. For $A\in\mathcal A_{T_i}$ and every
    natural $l$,
    \begin{align*}
        \lVert\tau_t(A)-E_{N_l(a_i)}(\tau_t(A))\rVert
            \leq\lVert A\rVert\min\Bigl(2,\;2e^RK_gC_k
                e^{2v_Rb_0e^{2R}|t|-l/2}\Bigr),
        \qquad C_k=\frac{2^kk!\,e^{1/2}}{1-e^{-1/2}}.
    \end{align*}
    For an arbitrary region $S$ with $A\in\mathcal A_S$ and its neighbourhood
    $N_l(S)=\bigcup_{x\in S}N_l(x)$, the error is at most
    $(2+4|b_0|e^{2R})|S|^2\lVert A\rVert K_gC_k
    e^{(2v_R|b_0|e^{2R}+1)|t|-l/2}$.
    For a single support this is~\cite[Lemma~4.1]{OpenAI2026AreaLaw}.
\end{theorem}
\begin{proof}\leanok
    \uses{thm:sites_region_expectation_localization, thm:sites_region_expectation,
        thm:dynamics_lieb_robinson, thm:dynamics_infinite_distance}
    The bound $2\lVert A\rVert$ holds since $E$ is a contraction and $\tau_t$ an
    isometry. For the other bound, Theorem~\ref{thm:sites_region_expectation_localization}
    reduces the error to the sum over sites $y$ outside the ball of the
    commutators with unitaries at $y$, which
    Theorems~\ref{thm:dynamics_lieb_robinson}
    and~\ref{thm:dynamics_infinite_distance} bound by
    $2e^R\lVert A\rVert e^{v|t|-d_G(a_i,y)}$. Grouping the sites by their
    distance and using the growth bound,
    $\sum_{d>l}(d+1)^ke^{-d}\leq C_ke^{-l/2}$.
\end{proof}
